\documentclass[tikz,border=4pt]{standalone}
% ==========================================================================
% tikz-tcee.tex — Estilos comunes de los gráficos TikZ del temario
% Autor: Víctor Gutiérrez Marcos
%
% Cada gráfico es un fichero standalone en fuentes/<tema>/graficos/:
%
%   \documentclass[tikz,border=4pt]{standalone}
%   % ==========================================================================
% tikz-tcee.tex — Estilos comunes de los gráficos TikZ del temario
% Autor: Víctor Gutiérrez Marcos
%
% Cada gráfico es un fichero standalone en fuentes/<tema>/graficos/:
%
%   \documentclass[tikz,border=4pt]{standalone}
%   % ==========================================================================
% tikz-tcee.tex — Estilos comunes de los gráficos TikZ del temario
% Autor: Víctor Gutiérrez Marcos
%
% Cada gráfico es un fichero standalone en fuentes/<tema>/graficos/:
%
%   \documentclass[tikz,border=4pt]{standalone}
%   \input{../../../latex/tikz-tcee}
%   \begin{document}
%   \begin{tikzpicture}[tcee]
%     \ejes{Q}{P}{6}{5}
%     \draw[curva1] (0.5,4.5) -- (5,0.8) node[etiqueta,right] {$D$};
%     \capa{2}{\draw[curva2] ...;}   % en el vídeo aparece en el paso 2
%   \end{tikzpicture}
%   \end{document}
%
% Paleta: la de la web (morado, dorado, salvia) más tres colores de apoyo
% distinguibles en blanco y negro por el trazo.
% ==========================================================================
\usepackage[T1]{fontenc}
\usepackage{newpxtext,newpxmath}
\usepackage{xcolor}
\usepackage{tikz,pgfplots}
\pgfplotsset{compat=1.18}
\usetikzlibrary{arrows.meta,calc,positioning,patterns,decorations.pathreplacing,intersections,shapes.misc,backgrounds}

\definecolor{tceemorado}{HTML}{5F2987}
\definecolor{tceedorado}{HTML}{B8860B}
\definecolor{tceeverde}{HTML}{3C7D3A}
\definecolor{tceeazul}{HTML}{1F5F99}
\definecolor{tceerojo}{HTML}{B22222}
\definecolor{tceegris}{HTML}{767171}
\definecolor{tceesalvia}{HTML}{E2EFD9}

% Capas para el vídeo: \capa{n}{…} dibuja su contenido solo a partir del paso n.
% En el PDF, la web y Word se ve todo; generar-video.py compila el gráfico una
% vez por paso con \def\tceepaso{k} para que cada elemento aparezca al narrarlo.
\providecommand{\tceepaso}{1000}
\newcommand{\capa}[2]{\ifnum#1>\tceepaso\relax\else#2\fi}

\tikzset{
  tcee/.style={line cap=round,line join=round,font=\small,>={Stealth[length=2.2mm]}},
  eje/.style={->,thick,black},
  curva1/.style={very thick,tceemorado},
  curva2/.style={very thick,tceeazul},
  curva3/.style={very thick,tceeverde},
  curva4/.style={very thick,tceedorado},
  curvanueva/.style={very thick,tceerojo},
  desplazada/.style={very thick,dashed},
  guia/.style={thin,dashed,tceegris},
  punto/.style={circle,fill=black,inner sep=1.3pt},
  etiqueta/.style={font=\small,inner sep=2pt},
  flecha/.style={->,thick,tceerojo},
  area/.style={fill=tceemorado,fill opacity=0.18,draw=none},
  area2/.style={fill=tceedorado,fill opacity=0.22,draw=none},
}

% \ejes{etiqueta x}{etiqueta y}{largo x}{alto y}
\newcommand{\ejes}[4]{%
  \draw[eje] (0,0) -- (#3,0) node[below left=2pt and 0pt] {#1};
  \draw[eje] (0,0) -- (0,#4) node[left] {#2};
  \node[below left] at (0,0) {$0$};}
% \proyeccion{(x,y)}{etq x}{etq y}: líneas guía a los ejes con etiquetas
\newcommand{\proyeccion}[3]{%
  \draw[guia] let \p1=#1 in (\x1,0) node[below] {#2} |- (0,\y1) node[left] {#3};}

%   \begin{document}
%   \begin{tikzpicture}[tcee]
%     \ejes{Q}{P}{6}{5}
%     \draw[curva1] (0.5,4.5) -- (5,0.8) node[etiqueta,right] {$D$};
%     \capa{2}{\draw[curva2] ...;}   % en el vídeo aparece en el paso 2
%   \end{tikzpicture}
%   \end{document}
%
% Paleta: la de la web (morado, dorado, salvia) más tres colores de apoyo
% distinguibles en blanco y negro por el trazo.
% ==========================================================================
\usepackage[T1]{fontenc}
\usepackage{newpxtext,newpxmath}
\usepackage{xcolor}
\usepackage{tikz,pgfplots}
\pgfplotsset{compat=1.18}
\usetikzlibrary{arrows.meta,calc,positioning,patterns,decorations.pathreplacing,intersections,shapes.misc,backgrounds}

\definecolor{tceemorado}{HTML}{5F2987}
\definecolor{tceedorado}{HTML}{B8860B}
\definecolor{tceeverde}{HTML}{3C7D3A}
\definecolor{tceeazul}{HTML}{1F5F99}
\definecolor{tceerojo}{HTML}{B22222}
\definecolor{tceegris}{HTML}{767171}
\definecolor{tceesalvia}{HTML}{E2EFD9}

% Capas para el vídeo: \capa{n}{…} dibuja su contenido solo a partir del paso n.
% En el PDF, la web y Word se ve todo; generar-video.py compila el gráfico una
% vez por paso con \def\tceepaso{k} para que cada elemento aparezca al narrarlo.
\providecommand{\tceepaso}{1000}
\newcommand{\capa}[2]{\ifnum#1>\tceepaso\relax\else#2\fi}

\tikzset{
  tcee/.style={line cap=round,line join=round,font=\small,>={Stealth[length=2.2mm]}},
  eje/.style={->,thick,black},
  curva1/.style={very thick,tceemorado},
  curva2/.style={very thick,tceeazul},
  curva3/.style={very thick,tceeverde},
  curva4/.style={very thick,tceedorado},
  curvanueva/.style={very thick,tceerojo},
  desplazada/.style={very thick,dashed},
  guia/.style={thin,dashed,tceegris},
  punto/.style={circle,fill=black,inner sep=1.3pt},
  etiqueta/.style={font=\small,inner sep=2pt},
  flecha/.style={->,thick,tceerojo},
  area/.style={fill=tceemorado,fill opacity=0.18,draw=none},
  area2/.style={fill=tceedorado,fill opacity=0.22,draw=none},
}

% \ejes{etiqueta x}{etiqueta y}{largo x}{alto y}
\newcommand{\ejes}[4]{%
  \draw[eje] (0,0) -- (#3,0) node[below left=2pt and 0pt] {#1};
  \draw[eje] (0,0) -- (0,#4) node[left] {#2};
  \node[below left] at (0,0) {$0$};}
% \proyeccion{(x,y)}{etq x}{etq y}: líneas guía a los ejes con etiquetas
\newcommand{\proyeccion}[3]{%
  \draw[guia] let \p1=#1 in (\x1,0) node[below] {#2} |- (0,\y1) node[left] {#3};}

%   \begin{document}
%   \begin{tikzpicture}[tcee]
%     \ejes{Q}{P}{6}{5}
%     \draw[curva1] (0.5,4.5) -- (5,0.8) node[etiqueta,right] {$D$};
%     \capa{2}{\draw[curva2] ...;}   % en el vídeo aparece en el paso 2
%   \end{tikzpicture}
%   \end{document}
%
% Paleta: la de la web (morado, dorado, salvia) más tres colores de apoyo
% distinguibles en blanco y negro por el trazo.
% ==========================================================================
\usepackage[T1]{fontenc}
\usepackage{newpxtext,newpxmath}
\usepackage{xcolor}
\usepackage{tikz,pgfplots}
\pgfplotsset{compat=1.18}
\usetikzlibrary{arrows.meta,calc,positioning,patterns,decorations.pathreplacing,intersections,shapes.misc,backgrounds}

\definecolor{tceemorado}{HTML}{5F2987}
\definecolor{tceedorado}{HTML}{B8860B}
\definecolor{tceeverde}{HTML}{3C7D3A}
\definecolor{tceeazul}{HTML}{1F5F99}
\definecolor{tceerojo}{HTML}{B22222}
\definecolor{tceegris}{HTML}{767171}
\definecolor{tceesalvia}{HTML}{E2EFD9}

% Capas para el vídeo: \capa{n}{…} dibuja su contenido solo a partir del paso n.
% En el PDF, la web y Word se ve todo; generar-video.py compila el gráfico una
% vez por paso con \def\tceepaso{k} para que cada elemento aparezca al narrarlo.
\providecommand{\tceepaso}{1000}
\newcommand{\capa}[2]{\ifnum#1>\tceepaso\relax\else#2\fi}

\tikzset{
  tcee/.style={line cap=round,line join=round,font=\small,>={Stealth[length=2.2mm]}},
  eje/.style={->,thick,black},
  curva1/.style={very thick,tceemorado},
  curva2/.style={very thick,tceeazul},
  curva3/.style={very thick,tceeverde},
  curva4/.style={very thick,tceedorado},
  curvanueva/.style={very thick,tceerojo},
  desplazada/.style={very thick,dashed},
  guia/.style={thin,dashed,tceegris},
  punto/.style={circle,fill=black,inner sep=1.3pt},
  etiqueta/.style={font=\small,inner sep=2pt},
  flecha/.style={->,thick,tceerojo},
  area/.style={fill=tceemorado,fill opacity=0.18,draw=none},
  area2/.style={fill=tceedorado,fill opacity=0.22,draw=none},
}

% \ejes{etiqueta x}{etiqueta y}{largo x}{alto y}
\newcommand{\ejes}[4]{%
  \draw[eje] (0,0) -- (#3,0) node[below left=2pt and 0pt] {#1};
  \draw[eje] (0,0) -- (0,#4) node[left] {#2};
  \node[below left] at (0,0) {$0$};}
% \proyeccion{(x,y)}{etq x}{etq y}: líneas guía a los ejes con etiquetas
\newcommand{\proyeccion}[3]{%
  \draw[guia] let \p1=#1 in (\x1,0) node[below] {#2} |- (0,\y1) node[left] {#3};}

\begin{document}
\begin{tikzpicture}[tcee]
  % Panel izquierdo: solución interior.
  % Recta x2 = 3 - 0,75 x1; U = x1 x2, tangente en x* = (2; 1,5) con U* = 3.
  \ejes{$x_1$}{$x_2$}{5}{4}
  \capa{1}{
    \fill[area] (0,0) -- (0,3) -- (4,0) -- cycle;
    \draw[curva1] (0,3) -- (4,0);
  }
  \capa{2}{
    \draw[curva2] plot[domain=0.45:4.6,samples=80,smooth] (\x,{1.6/\x}) node[etiqueta,right] {$U^0$};
    \draw[curva3] plot[domain=0.8:4.6,samples=80,smooth] (\x,{3/\x}) node[etiqueta,right] {$U^*$};
    \draw[curva4] plot[domain=1.35:4.6,samples=80,smooth] (\x,{5/\x}) node[etiqueta,right] {$U^2$};
    \node[punto] (E) at (2,1.5) {};
    \node[etiqueta,above right] at (E) {$x^*$};
    \proyeccion{(E)}{$x_1^*$}{$x_2^*$}
    \node[etiqueta,anchor=west] at (2.6,2.6) {$RMS_{12}=p_1/p_2$};
    \draw[guia,solid] (2.7,2.45) -- (2.1,1.6);
  }
  % Panel derecho: solución de esquina en el eje x1.
  % Recta x2 = 2 - 0,5 x1 (p1/p2 = 0,5); U = x1 + ln(1+x2): x2 = e^{c-x1} - 1,
  % con RMS = 1 + x2 >= 1 > p1/p2, así que la curva más alta alcanzable (c = 4)
  % toca la recta solo en (4, 0).
  \begin{scope}[xshift=6.8cm]
    \ejes{$x_1$}{$x_2$}{5.6}{4}
    \capa{3}{
      \fill[area] (0,0) -- (0,2) -- (4,0) -- cycle;
      \draw[curva1] (0,2) -- (4,0);
      \draw[curva2] plot[domain=1.7:3.2,samples=60,smooth] (\x,{exp(3.2-\x)-1});
      \draw[curva3] plot[domain=2.5:4,samples=60,smooth] (\x,{exp(4-\x)-1});
      \draw[curva4] plot[domain=3.3:4.8,samples=60,smooth] (\x,{exp(4.8-\x)-1});
      \node[etiqueta,above] at (1.7,{exp(1.5)-1}) {$U^0$};
      \node[etiqueta,above] at (2.5,{exp(1.5)-1}) {$U^*$};
      \node[etiqueta,above] at (3.3,{exp(1.5)-1}) {$U^2$};
      \node[punto] at (4,0) {};
      \node[below] at (4,0) {$x_1^*=\overline{W}/p_1$};
      \node[etiqueta,above right] at (4,0) {$x^*$};
      \node[etiqueta,align=left,anchor=west] at (3.9,2.1) {$x_2^*=0$\\$RMS_{12}\geq p_1/p_2$};
    }
  \end{scope}
\end{tikzpicture}
\end{document}
